% ----------------------------------------------------------------------------
% Vorlage Abschlussarbeit Informatik THM (minimal)
%
% Copyright (c) 2016 - 2020 by Burkhardt Renz. All rights reserved.
% Die Vorlage für eine Abschlussarbeit in der Informatik am Fachbereich
% MNI der THM ist lizenziert unter einer Creative Commons
% Namensnennung-Nicht kommerziell 4.0 International Lizenz.
%
% $Id: vorlage.tex 979 2020-08-24 07:03:19Z br $
% ----------------------------------------------------------------------------

\documentclass[%
	BCOR=8.25mm,         % Bindekorrektur
	DIV=12,              % Satzspiegel
	parskip=half,				 % Abstand zwischen Absätzen
	bibliography=totoc,	 % Literaturverzeichnis im Inhaltsverzeichnis
	headsepline=on,      % Trennlinie Kolumnentitel
	]{scrbook}

%% Präambel
\usepackage[english, ngerman]{babel} % deutsche typogr. Regeln + Trenntabelle
\usepackage[T1]{fontenc}             % interner TeX-Font-Codierung
\usepackage{helvet}                 % Font Latin Modern
\usepackage[utf8]{inputenc}          % Font-Codierung der Eingabedatei
\usepackage[babel]{csquotes}         % Anführungszeichen
\usepackage{graphicx}                % Graphiken
\usepackage{booktabs}                % Tabellen schöner
\usepackage{listingsutf8}            % Listings mit Einstellungen
\lstset{basicstyle=\small\ttfamily,
	tabsize=2,
	basewidth={0.5em,0.45em},
	extendedchars=true}
\lstset{literate=%                   % Umlaute in Listings
  {Ö}{{\"O}}1
  {Ä}{{\"A}}1
  {Ü}{{\"U}}1
  {ß}{{\ss}}2
  {ü}{{\"u}}1
  {ä}{{\"a}}1
  {ö}{{\"o}}1
	{~}{{\textasciitilde}}1} 
\usepackage{amsmath}	               % Mathematik
\usepackage[pdftex]{hyperref}       
\hypersetup{
	bookmarksopen=true,
	bookmarksopenlevel=3,
	colorlinks,
	citecolor=blue,
	linkcolor=blue,
}
\usepackage{scrhack}				% unterdrückt Fehlermeldung von listings
\usepackage[
	sorting = none,
	style = authoryear,
	backend = biber]{biblatex}   % Biblatex für das Literaturverzeichnis
\addbibresource{Hadoop.bib}			% Import der Bibliothek von Zotero

%% Nummerierungstiefen
\setcounter{tocdepth}{3}             % 3 Stufen im Inhaltsverzeichnis
\setcounter{secnumdepth}{3} 		     % 3 Stufen in Abschnittnummerierung

% ----------------------------------------------------------------------------
\begin{document}

\frontmatter

%% Titelseite
\begin{titlepage}
	\begin{center}
	\includegraphics[width=0.9\textwidth]{img/mni-logo}
	
	\vspace{5cm}	

	\vspace{1cm}	

	\huge\textbf{\sffamily Apache Hadoop}

	\normalsize
	\vspace{1cm}	

	Hauptseminar "Cloud-Plattformen und Big Data" \\
	Dozent Steffen Rupp

	von \\[1cm]	

	\textbf{René Gentzen}\\ [.5cm] 
	rene.gentzen@mni.thm.de\\ [.5cm] 
	im WS22/23
	\end{center}
	\vfill
\end{titlepage}

%% Verzeichnissse
\tableofcontents

\listoffigures
\listoftables
\lstlistoflistings

\mainmatter 
\pagestyle{headings}
\chapter{Hadoop Grundlagen}

``Die Apache Hadoop Softwarebibliothek ist ein Framework, das die über Computercluster verteilte Verarbeitung großer Datensätze mit einfachen Programmiermodellen ermöglicht. Es ist so konzipiert, dass es von einzelnen Servern bis hin zu Tausenden von Rechnern skaliert werden kann, von denen jeder lokale Rechenleistung und Speicherplatz bietet. Anstatt sich auf Hardware zu verlassen, um eine hohe Verfügbarkeit zu gewährleisten, ist die Bibliothek selbst so konzipiert, dass sie Ausfälle auf der Anwendungsebene erkennt und bewältigt, so dass ein hochverfügbarer Dienst auf einem Cluster von Computern bereitgestellt wird, von denen jeder für sich für Ausfälle anfällig sein kann.''\cite{noauthor_apache_nodate}\\
So beschreibt (übersetzt aus dem Englischen) die Apache Software Foundation ihr Top Level Projekt Apache™ Hadoop®.  Diese Arbeit wird einen pragmatischen Überblick über Hadoop und die Komponenten im Hadoop Ecosystem geben. Dabei soll der Fokus nicht auf technischen Details, sondern auf Anwendungsorientiertheit leigen. Es sollen konkrete Anwendungsfälle zu den einzelnen Komponenten besprochen werden und nicht abstrakte Architekturkonzepte.

\section{Hadoop Historie}
\subsection{Anforderungen von Big Data}
``Der Begriff „Big Data“ bezieht sich auf Datenbestände, die so groß, schnelllebig oder komplex sind, dass sie sich mit herkömmlichen Methoden nicht oder nur schwer verarbeiten lassen.''\cite{noauthor_big_nodate} \\ 
Schon Anfang der Neunziger war es nicht mehr praktikabel, Webseiten händisch, zum Beispiel in ''Web Directories'', zu katalogisieren. Man wollte Nutzern trotzdem die Möglichkeit geben, Informationen durch das Durchsuchen zentraler Anlaufstellen ausfindig zu machen. Automatisierte Tools, die sogenannten ''Web Crawler'' wurden erfunden, um diese Arbeit zu übernehmen.\cite{griffiths_search_2007} \\ 
Das Internet erlebte in den letzten Jahren des 20. Jahrhunderts ein explosionsartiges Wachstum an Nutzern und Webseiten, und damit auch an Informationen, die katalogisiert werden mussten.\cite{zakon_hobbes_2018} 
Um eine immer größer werdende Menge an Informationen verarbeiten zu können, gibt es zwei Ansätze der Skalierung: Vertikale und horizontale Skalierung. Diese sollen in den folgenden Abschnitten erläutert werden.

\subsection{Vertikale Skalierung}
Bei der vertikalen Skalierung (''scaling up'') werden \emph{einem} System mehr Ressourcen wie zum Beispiel größerer Speicher, oder eine schnellere CPU hinzugefügt. Dadurch bekommt man einen Performance-Gewinn: Man kann mehr Daten speichern, oder Berechnungen werden schneller fertig gestellt.
Ein großer Vorteil der vertikalen Skalierung ist, dass Anwendungsprogramme in der Regel nicht angepasst werden müssen, um vom diesem Performance-Wachstum zu profitieren. Wenn man eine 5TB große Festplatte gegen eine 10TB Festplatte austauscht, dann hat man den Speicherplatz eines Servers vertikal skaliert. Die darauf laufenden Programme müssen nicht angepasst werden, sondern man kann einfach doppelt so viele Daten speichern.\cite{beaumont_how_2014}\\
Vertikale Skalierung hat drei große Nachteile: Erstens kann man nicht unbegrenzt vertikal skalieren. Ein Server kann physisch nur eine begrenzte Anzahl an Hardware aufnehmen. Zweitens wächst die Performance eines Systems bei vertikaler Skalierung höchstens linear, die Kosten allerdings nicht. Heutzutage kann man gerade bei Cloud-Anbietern sehr leistungsfähige Systeme bei linearem Preisanstieg mieten.\cite{noauthor_pricing_nodate} Sucht man aber noch mehr Performance in \emph{einem} System, dann steigen die Kosten exponentiell\cite{athow_at_2020}. Drittens skalieren nicht alle Faktoren in einem System gleich gut vertikal. Die Speicherkapazität von SSDs ist zum Beispiel seit 1978 von 45MB auf 100TB gestiegen (Faktor $2222,22 \cdot 10^{3}$), während sich die Datenrate nur von 1.5MB/s auf 500/460MB/s (Sequential Read/Write) erhöht hat (Faktor $0,333 \cdot 10^{3}$).\cite{noauthor_who_nodate}\cite{athow_at_2020}  


\subsection{Horizontale Skalierung}
Bei der horizontalen Skalierung (''scaling out'')

\subsection{GFS und MapReduce}
\section{Hadoop 1.0 Aufbau}
\subsection{HDFS in depth theoretisch}
\subsection{MapReduce in depth theoretisch}

\section{Hadoop 2.x mit YARN}
\subsection{Neue Möglichkeiten mit YARN}
\subsection{Neuer Aufbau, neue Verpflichtungen}
\subsection{Code einer YARN Application}

\section{Hadoop 3.x anreißen}
\chapter{Das Hadoop Ecosystem}
\section{Storage}
\subsection{HDFS}
\subsection{HBase}

\section{Management / Configuration}
\subsection{YARN}
\subsection{Oozie}
\subsection{ZooKeper}
\subsection{Ambari}

\section{Datentransfer}
\subsection{Sqoop}
\subsection{Kafka}
\subsection{AVRO}

\section{Processing}
\subsection{MapReduce}
\subsection{Pig}
\subsection{Hive}
\subsection{Flume}
\subsection{Mahout}
\subsection{Spark}
\subsection{Solr}

\chapter{Hadoop Operations}
\section{Hadoop Setup} 
\subsection{Single Node Setup}
\subsection{Pseudo-Distributed Setup} 
\subsection{Fully-distributed Cluster}
\subsection{Docker Images}
\subsection{VM Distributionen}
\subsection{Hadoop in der Cloud}
HDFS oder Cloud FS 

\section{Praxis}
\subsection{Start eines Single Node Clusters lokal}
\subsection{Erste Übung zum Umgang mit dem HDFS}
\subsection{MapReduce Workflow: Erstellung eines MapReduce Jobs, Kopieren auf den Name Node und Ausführung}
\subsection{YARN Application auf einem Cluster in der Cloud laufen lassen}
\chapter{ETL mit Pig}
Auch wenn es vermehrt von Spark verdrängt wird
\section{Anwendungsfälle}
Welche neuen Dinge ermöglicht dieses Tool
Eine Zeile Pig Latin entspricht vielen Zeilen MapReduce
\section{Architektur}
\section{Pig Latin}
\section{Praxis}
\subsection{Hinzufügen zum Cluster}
\subsection{Anwendung auf dem Cluster}
\chapter{Data Ingestion}
\section{Sqoop}
\section{Flume}
\chapter{Datawarehousing mit Hive}
\section{Anwendungsfälle}
Welche neuen Dinge ermöglicht dieses Tool
\subsection{Unterschiede zu Pig}
\section{Architektur}
\section{Interaktion}
\subsection{HiveQL}
\subsection{CLI}
\subsection{Java API}
\section{Praxis}
\subsection{Hinzufügen zum Cluster}
\subsection{Einrichtung einer Datenbank}
\subsection{Einlesen von Daten im CLI}
\subsection{Einlesen von Daten mit Sqoop}
\subsection{Absetzen einer Query}
\chapter{NoSQL mit HBase}
\section{Anwendungsfälle}
Welche neuen Dinge ermöglicht dieses Tool
\subsection{CAP-Theorem}
\subsection{ACID und BASE}
\section{Architektur}
\section{Interaktion}
\subsection{HBase Shell}
\subsection{Java API}
\section{Praxis}
\subsection{Hinzufügen zum Cluster}
\subsection{Einrichtung einer Datenbank}
\subsection{Einlesen von Daten}
\subsection{Datenmigration aus einem RDBMS}
\subsection{Absetzen einer Query}
\chapter{Streaming mit Kafka}
\section{Anwendungsfälle}
Welche neuen Dinge ermöglicht dieses Tool
Ersetzt durch Spark Streaming
\section{Architektur}
\section{Interaktion}
\subsection{Who knows}
\section{Praxis}
\subsection{Hinzufügen zum Cluster}
\subsection{Maybe, vielleicht kann man ja was zeigen}
\chapter{Hadoop heute}
\section{Aktuelle Anwendungsbeispiele zu Hadoop}
\subsection{AirBnB}
\section{Apache Spark als Gold Standard}
\subsection{Kann eh alles besser}

\backmatter 

\appendix

\printbibliography

\end{document}
% ----------------------------------------------------------------------------
